\section{Finite-Step Right-Generator Certificate}
\label{sec:finite-step}

An asymptotic statement $O(t^5)$ is insufficient at the target tolerance.
The certified step count is obtained from a finite-$t$ right-generator bound,
including explicit subleading terms and a rational tail.

\subsection{From logarithmic defect to right generator}

Let $S(t)$ denote one formula step and define its right generator
\begin{equation}
  K(t)=S'(t)S(t)^{-1}.
  \label{eq:right-generator}
\end{equation}
For $Z(t)=\log S(t)$, differentiation of the exponential gives
\begin{equation}
  K(t)=\operatorname{dexp}_{Z(t)}(Z'(t))
      =\sum_{n\geq0}\frac{1}{(n+1)!}\ad_{Z(t)}^n Z'(t).
  \label{eq:dexp}
\end{equation}
Substituting \cref{eq:symmetric-log} and collecting powers yields
\begin{equation}
  K(t)-A=t^4D_4+t^5D_5+t^6D_6+t^7D_7+R_{\geq8}(t),
  \label{eq:right-ledger}
\end{equation}
where, writing $B=E_5$ and $C=E_7$, the exact coefficients through degree
seven are
\begin{align}
 D_4&=5B, &
 D_5&=2[A,B],\\
 D_6&=7C+\tfrac23[A,[A,B]], &
 D_7&=3[A,C]+\tfrac16[A,[A,[A,B]]].
\end{align}
The implementation derives these coefficients from the truncated series
rather than trusting the printed display.  The identity clarifies why an
exact leading logarithmic coefficient is not yet a global error certificate:
commutators generated by the differential of the exponential contribute at
the next powers.

\subsection{Duhamel conversion to global error}

Let $U(t)=e^{tA}$ in anti-Hermitian convention.  Variation of constants gives
the one-step estimate
\begin{equation}
 \opnorm{S(\delta)-U(\delta)}
 \leq \int_0^\delta \opnorm{K(s)-A}\,ds.
 \label{eq:one-step-duhamel}
\end{equation}
Unitary telescoping over $r$ equal steps, with $\delta=T/r$, then yields
\begin{equation}
 \opnorm{S(T/r)^r-U(T)}
 \leq r\int_0^{T/r}\opnorm{K(s)-A}\,ds.
 \label{eq:global-duhamel}
\end{equation}
Combining \cref{eq:right-ledger,eq:global-duhamel}, a certificate with norm
upper bounds $b_j\geq\opnorm{D_j}$ and a nonnegative tail majorant $\rho(s)$
establishes
\begin{equation}
 \mathcal E(r)=
 r\left(\sum_{j=4}^{7}\frac{b_j}{j+1}
                  \left(\frac{T}{r}\right)^{j+1}
 +\int_0^{T/r}\rho(s)\,ds\right).
 \label{eq:finite-step-bound}
\end{equation}
All quantities in the machine certificate are rationals.  Decimal values in
this paper are outward-rounded summaries only.

\subsection{Rational tail closure}

The exact ledger is stopped after D7.  Higher degrees are bounded by a
geometric majorant derived from the absolute stage coefficients and local
commutator growth.  In abstract form, if the degree-$k$ coefficient is at most
$M q^{k-k_0}$ for $0\leq q s<1$, then
\begin{equation}
 \rho(s)\leq \frac{M s^{k_0}}{1-qs}.
 \label{eq:geometric-tail}
\end{equation}
The verifier checks the rational preconditions, including positivity of the
denominator on $[0,T/r]$, and integrates an outward upper enclosure.  This is
deliberately more conservative than extrapolating the observed low-degree
coefficients.  It ensures that no uncomputed term is silently discarded.

\begin{theorem}[Soundness of an accepted finite-step certificate]
\label{thm:soundness}
Suppose the verifier accepts (i) interval enclosures for D4--D7, (ii) a valid
operator-norm certificate for each declared coefficient, (iii) a tail
majorant satisfying \cref{eq:geometric-tail} on $[0,T/r]$, and (iv) exact
rational evaluation $\mathcal E(r)\leq\eps$.  Then
\begin{equation}
  \opnorm{S(T/r)^r-e^{TA}}\leq\eps.
\end{equation}
\end{theorem}

\begin{proof}
The accepted coefficient and tail checks bound
$\opnorm{K(s)-A}$ term by term on the complete integration interval.
Integrating preserves the inequality.  Equation~\eqref{eq:global-duhamel}
then bounds the global error, and the exact rational comparison with $\eps$
completes the proof.
\end{proof}

\subsection{Adjacent-step minimality under the certified bound}

The search evaluates integers, not a rounded real-valued estimate.  It returns
the least $r$ accepted by \cref{eq:finite-step-bound} and emits the rejected
neighbor.  For this benchmark,
\begin{align}
 \mathcal E(97)&\leq9.958938494314325\times10^{-7}<10^{-6},\\
 \mathcal E(96)&\leq1.050565873970784\times10^{-6}>10^{-6}.
\end{align}
The second inequality says that 96 is not certified by \emph{this bound}; it
does not prove that the physical simulation error at 96 steps exceeds the
tolerance.  This distinction is retained throughout the paper.
